\section{Introduction}
This document describes ProcControlAPI, an API and library for controlling
processes. ProcControlAPI runs as part of a controller process
and manages one or more target processes. It allows the controller process
to perform operations on target
processes, such as writing to memory, stopping and running threads, or
receiving notification when certain events occur. ProcControlAPI presents
these operations through a platform-independent API and high-level
abstractions. Users can describe what they want ProcControlAPI to do, and
ProcControlAPI handles the details.

An example use for ProcControlAPI would be as the underlying mechanism for a
debugger. A user writing a debugger could provide their own user interface and
debugging strategies, while using ProcControlAPI to perform operations such as
creating processes, running threads, and handling breakpoints.

The interface for ProcControlAPI can be generally divided into two parts: an
interface for managing a process (e.g., reading and writing to target process
memory, stopping and running threads), and an interface for monitoring a target
process for certain events (e.g., watching the target process for fork or
thread creation events). The manager interface uses a set of C++ objects to
represent a target process and its threads, libraries, registers and other
interesting aspects. Operations performed on these C++ objects in the
controller process are translated into corresponding operations on the target
process. The event interface uses a callback system to notify the
ProcControlAPI user of interesting events in the target process. 

\subsection{Simple Example}

As an example, consider the code in Figure 1 that creates a target process and
prints a message whenever that target process creates a new thread. Details
on the API function used in this example can be found in latter sections of
this manual, but we will provide a high level description of the operations
here. Note that proper error handling and checking have been left out for
brevity.

\begin{lstlisting}[language=c++]
  #include "Event.h"
  #include "PCProcess.h"

  #include <iostream>
  #include <string>

  namespace pc = Dyninst::ProcControlAPI;

  pc::Process::cb_ret_t on_thread_create(pc::Event::const_ptr ev) {
    // Callback when the target process creates a thread.
    auto new_thrd_ev = ev->getEventNewThread();
    auto new_thrd = new_thrd_ev->getNewThread();

    std::cout << "Got a new thread with LWP " << new_thrd->getLWP() << std::endl;
    return Process::cbDefault;
  }

  int main(int argc, char *argv[]) {
    std::vector<std::string> args;

    // Create a new target process
    std::string exec = argv[1];
    for (int i = 1; i < argc; i++) {
      args.emplace_back(argv[i]);
    }

    pc::Process::ptr proc = pc::Process::createProcess(exec, args);

    // Tell ProcControlAPI about our callback function
    pc::Process::registerEventCallback(pc::EventType::ThreadCreate,
                                       on_thread_create);

    // Run the process and wait for it to terminate.
    proc->continueProc();

    while (!proc->isTerminated()) {
      pc::Process::handleEvents(true);
    }
  }
\end{lstlisting}

We ask ProcControlAPI to create a new Process using the
given arguments. ProcControlAPI will spawn a new target process
 and leave it in a stopped state to prevent it from executing.

\begin{enumerate}
\item After creating the new target process we register a callback
function. We ask ProcControlAPI to call our
function, on\_thread\_create, when an event of type EventType::ThreadCreate
occurs in the target process. 

\item The on\_thread\_create function takes a pointer to an object of type
Event and returns a Process::cb\_ret\_t. The Event describes the target
process event that triggered
this callback. In this case, it provides information about the new thread in
the target process. It is worth noting that Event::const\_ptr is a not a
regular pointer, but a reference counted shared pointer. This means that we
do not have to be concerned with cleaning the Event- it
will be automatically cleaned when the last reference disappears. The
Process::cb\_ret\_t describes what action should be taken on the process in
response to this event, which is described in more detail in section 6. 

\item The Event class has several child classes, one of which is EventNewThread.
We start by casting the Event into an EventNewThread and then extract
information about the new thread from the EventNewThread.

\item In step 6, we\'ve finished handling the new thread
event and need to tell ProcControlAPI what to do in response to this event.
For example, we could choose to stop the process from further execution by
returning a value of Process::cbProcStop. Instead, we choose let
ProcControlAPI take its default action for an EventNewThread by returning
Process::cbDefault, which is to continue the process and its new thread (which
were both stopped before delivery of the callback).

\item The registering of our callback in step 3 did not actually trigger any
calls to the callback function- the target process was
created in a stopped state and has not yet been able to create any threads.
We tell ProcControlAPI to continue the target process in this step, which
allows it to execute and possibly start generating new events. 

\item In this step we wait for the target process to finish executing and
terminate. Calling Process::handleEvents
blocks the controller process until an event occurs, allowing us to wait for
events without needing to spin the controller process on the CPU. 
\end{enumerate}
